\scanpage{29}
\deleted{\uline{Exemple} Prenons \unsure{$\varphi$ non dég|pour \illegible{1} dég}}
\[
\left|\begin{aligned}
\varphi:\bigwedge^2 V&\to k\qquad\text{non dégénérée}\\
&\inserted{\text{i.e. }\varphi\in\bigwedge^2\check V}\\
&\deleted{\unsure{\text{d'où }}\check\varphi\in\unsure{\bigwedge}}\\
\chi&=\varphi^2\in\bigwedge^4\check V\qquad \chi:\bigwedge^4 V\to k
\end{aligned}\right.
\]
\deleted{Alors $\chi$ comme forme sur $\bigwedge^2 V$ est \unsure{non}}

\uline{Ex} Prenons $\chi'$ non dégénérée, \deleted{\unsure{$\alpha$ un isom.}} $\alpha=\mathrm{id}$ i.e. $\widetilde W=\bigwedge^2 V$.

Alors un choix convenable de $\xi\in\unsure{\bar W|\widetilde W}=\bigwedge^2 V$ donne n'importe quelle $\varphi$ voulue, \unsure{au vecteur scalaire|au facteur scalaire|à un facteur scalaire}, donc n'importe quelle $\check\varphi$ \inserted{$\in\bigwedge^2 V$} non dégénérée, au \unsure{vecteur scalaire|facteur scalaire}. Considérons alors $\check\varphi^{\unsure{2|1|{}}}$ comme élément de \deleted{$\bigwedge^2 V\otimes$} l'algèbre $\bigwedge^{\unsure{*|\bullet|\vee}}V\otimes\bigwedge^\bullet V$ [on écrira $\bigwedge^2 V\subset V\otimes V$] et $\check\varphi^{2}$ comme élément de $\bigwedge^2 V\otimes\bigwedge^2 V\subset\bigwedge^*V\otimes\bigwedge^*V$ ; le carré de $\check\varphi$ dans $\bigwedge^*V$ \inserted{(qui est $\in\bigwedge^4 V$)} est l'image de $\check\varphi^{\unsure{(2)|(n)}}$ précédent par $\bigwedge^2 V\otimes\bigwedge^2 V\longrightarrow\bigwedge^4 V$. Utilisant la forme $Q\unsure{:|=}Q_{\chi'}$ définie par $\chi'$, on voit que $\check\varphi^{(2)}$ définit un endom de $\bigwedge^2 V$, \uline{soit $u$}. Je dis que l'orbite sous $u$ des \deleted{\unsure{él}}éléments $\xi$ et $\deleted{\unsure{d}}\check\varphi$ de $\bigwedge^2 V$ peuvent être infinis, si $\dim V\geq4$ (ou 6?) ($\unsure{W|V}$ est de dim paire) \deleted{\illegible{1}} et \inserted{les espaces engendrés par les uns et \deleted{\unsure{pl}} les autres} de dim non bornée (pour dim $V$ croissante\ldots).
% archive sheet number: 28
